\section{The QPASE Scheme}
\label{sec:qpase}

This section restates QPASE as Sections III and IV of the paper specify it~\cite{jiang2024qpase}.

\paragraph{Model.} The user remembers only $ID_u$ and $psw_u$ and stores no key. Each of $N$ semi-honest servers
holds a share $k_i$, a record $(ID_u, pk_u, \mu_u)$ per user, and encrypted files. A quantum adversary controls the
network and may corrupt fewer than $t$ servers. It tries offline and online password guessing and chosen-keyword
attacks on the index. The servers refuse a user after $\mu$ login attempts in one epoch.

\paragraph{Setup.} Setup publishes the ring $\Rq$ with $q = 2^{28} - 57$, the vectors $a, a_0, a_1$, $\sigma$,
$\mu$, $N$ and $t$. Table II of the paper gives two ring degrees: PARAM~I with $n = 512$ (92-bit quantum security)
and PARAM~II with $n = 595$ (128-bit), which its experiments use.

\begin{protocol}{Register and Login (paper Sections IV-B and IV-C)}
\begin{enumerate}[1.]
  \item \textbf{Register.} $U$ runs the TOPRF with every server to get $\Ku = F_K(x)$, derives
        $(pk_u, sk_u) \leftarrow \mathrm{Gen}(\Ku)$ and sends $pk_u$. Each $S_i$ stores $(ID_u, pk_u)$ and sets
        $\mu_u = 0$.
  \item \textbf{L1.} $U$ runs the TOPRF with $t$ servers. Each checks $\mu_u < \mu$, increments $\mu_u$, answers
        and returns $pk_u$.
  \item \textbf{L2.} $U$ recomputes $pk'_u$ from $\Ku$. If $pk'_u \neq pk_u$ the password was wrong. Otherwise it
        signs $\sigma_u = \mathrm{Sign}(sk'_u, (ID_u, sid))$.
  \item \textbf{L3.} $S_i$ verifies $\sigma_u$ and lets $U$ upload and search.
\end{enumerate}
\end{protocol}

\begin{protocol}{Outsource and Retrieve (paper Sections IV-D and IV-E)}
\begin{enumerate}[1.]
  \item \textbf{O1.} For a document $d$ with keyword $w$ and item id $i$, $U$ samples $\rho$ and computes
        $\dsk = \mathrm{HKDF}(\Ku, w)$, $v = \mathrm{PRF}(\dsk, \rho)$, $C = \mathrm{Enc}(\Ku, d)$,
        $\sigma_C = \mathrm{Sign}(sk_u, (\rho, v, C))$ and $Ct = \mathrm{Enc}(\dsk, (\rho, v, \sigma_C))$. It sends
        $(Ct, C, i)$ with a signature to one server.
  \item \textbf{O2.} The server verifies the signature and stores $(Ct, C, i)$.
  \item \textbf{R1.} $U$ sends $\dsk$ and $\mathrm{Sign}(sk_u, \dsk)$ to every server.
  \item \textbf{R2.} Each $S_i$ decrypts every stored $Ct$ with $\dsk$ and returns the items where
        $\mathrm{Ver}(pk_u, (\rho, v, C), \sigma_C)$ holds and $v = \mathrm{PRF}(\dsk, \rho)$.
  \item \textbf{R3.} $U$ repeats the check on each returned item and decrypts $d = \mathrm{Dec}(\Ku, C)$.
\end{enumerate}
\end{protocol}

The server finds a match by trying $\dsk$ on every item it stores, and the user sends the same $\dsk$ every time it
searches for $w$. \Cref{sec:counter-index} builds on both facts.

\paragraph{Server key update.} Each server deals shares of a random polynomial with constant term 0, with
commitments $H(\alpha_k)$ to its coefficients. Each $S_i$ checks $H(F_{ji}) = \sum_k i^k H(\alpha_k)$ and adds the
shares to $k_i$. The master key $K$ does not change, so users keep $\Ku$ and their files. The paper also sketches
multi-keyword queries (Section IV-G), which we left out.
